\section{The proposed test statistic}\label{sec:test}

\subsection{Grouped residuals and a basis of calibration shapes}

Sort the $n$ observations by fitted risk $\hat\pi_i$, ties in the order of the records, and cut them
into $G$ groups of equal frequency, as the Hosmer--Lemeshow test does. Write $O_g$ for the events observed in group $g$,
$E_g=\sum_{i\in g}\hat\pi_i$ for those expected, $V_g=\sum_{i\in g}\hat\pi_i(1-\hat\pi_i)$, $\pibar_g$
for the mean fitted risk of the group, and
\begin{equation}\label{eq:resid}
  r_g=\frac{O_g-E_g}{\sqrt{V_g}},\qquad r=(r_1,\dots,r_G)^\top .
\end{equation}
The Hosmer--Lemeshow statistic is $\lVert r\rVert^2$. It responds to a departure in any of the $G$
directions of residual space and spends a degree of freedom on each.

Suppose the fitted model is wrong through a misspecified link or an omitted smooth term, so that the
true event rate among patients of predicted risk $\pibar_g$ is $\pibar_g+\delta(\pibar_g)$ for a smooth
function $\delta$. Then $\mathbb{E}[r_g]\approx|g|\,\delta(\pibar_g)/\sqrt{V_g}$ is itself a smooth
function of $\pibar_g$, and a smooth function sampled at $G$ ordered points is described almost
entirely by a few shape components: a tilt, a bend and a second bend. The mean of $r$ therefore lies
close to a low-dimensional subspace, while the Hosmer--Lemeshow statistic weights all $G$ directions
equally.

Let $Z\in\mathbb{R}^{G\times d}$, with $d$ small, have columns that are smooth functions of the
calibration position. The test statistic is the squared length of the projection of $r$ onto the
span of $Z$,
\begin{equation}\label{eq:stat}
  S=r^\top Z(Z^\top Z)^{-1}Z^\top r=\lVert P_Z r\rVert^2 ,
  \qquad P_Z=Z(Z^\top Z)^{-1}Z^\top .
\end{equation}
The choice $Z=I_G$ returns the Hosmer--Lemeshow statistic, and the single column
$a_g=(1-2\pibar_g)/\sqrt{V_g}$ returns, up to scaling, the first-order correction of
\citet{osius1992normal} and \citet{farrington1996}. The groups and the residual vector are those of
the Hosmer--Lemeshow test; only the final combination differs. The weights act on groups, through the
basis, and never on the residual of an individual record, which is what Section~\ref{sec:mechanism}
turns on.

\subsection{The bases}\label{sec:bases}

Three bases are used, each a choice of alternative rather than a tuning parameter. With
$\bar\eta_g=\operatorname{logit}(\pibar_g)$:
\begin{description}
\item[\DEFp{3}, the default.] Orthogonal polynomials $P_1,P_2,P_3$ of $\pibar_g$, the tilt, bend and
double bend of the calibration curve; $d=3$.
\item[\DEFstk{}.] The grouped form of Stukel's constructed variables, $\bar\eta_g$,
$\bar\eta_g^2\mathbf 1\{\bar\eta_g\ge0\}$ and $-\bar\eta_g^2\mathbf 1\{\bar\eta_g<0\}$, which let the link
bend separately in each tail; $d=3$. It is a grouped version of Stukel's score test
\citep{stukel1988generalized}.
\item[\DEFsym{}.] The single column $\bar\eta_g|\bar\eta_g|$, the symmetric member of Stukel's
family, in which both tails bend together; $d=1$.
\end{description}
We write EDGE for the test in general and name the basis where it matters.

The columns of $Z$ may be used as they stand, the \emph{unit} form, or first scaled by $\sqrt{V_g}$,
the \emph{score} form, in which $Z^\top r$ is the score for adding the grouped shape to the model. The
score form moves weight towards the extreme risk groups. The unit form with the cubic basis is the
default and is used throughout unless stated. With the symmetric basis the score form is the better
choice when discrimination is high (an area under the curve near $0.9$ in our designs) and when a
covariate is strongly skewed, provided the data are clean: because it leans on the extreme groups, it
is the form most affected by corrupted records (Supporting Information S1).

\subsection{The null distribution}\label{sec:null}

The model is fitted to the same data, so $r$ is not a vector of independent standard normals. With
$X$ the model matrix, $W$ the fitted weights and $U$ the $G\times p$ matrix with rows
$U_g=\sum_{i\in g}\hat\pi_i(1-\hat\pi_i)x_i^\top/\sqrt{V_g}$,
\begin{equation}\label{eq:omega}
  r\ \dot\sim\ N(0,\Omega),\qquad \Omega=I_G-U(X^\top WX)^{-1}U^\top .
\end{equation}
Fitting the model absorbs part of each group's residual, so the diagonal of $\Omega$ is below one.
When a model is validated on new data with its coefficients frozen, nothing is estimated from those
data and $\Omega=I_G$. No score equation then absorbs the overall level, so a constant column is added
to the basis and the statistic has $d+1$ degrees of freedom; we call this the external mode. It uses
only the outcomes and the predicted probabilities, so it applies to the predictions of any model,
logistic or not, provided they were produced without the validation data.

\begin{theorem}[Null distribution]\label{thm:null}
Under the null hypothesis and conditions (A1)--(A4) of Appendix~\ref{app:proofs},
\begin{equation}\label{eq:wchi2}
  S\ \dot\sim\ \sum_{j=1}^{d}\lambda_j\,\chi^2_{1,j},\qquad
  \{\lambda_j\}=\text{eigenvalues of }(Z^\top Z)^{-1}Z^\top\Omega Z ,
\end{equation}
with $0\le\lambda_j\le1$ for every $j$.
\end{theorem}

The $p$-value is computed from the Satterthwaite scaled chi-squared approximation to this weighted
sum \citep{satterthwaite1946}, which is exact when $d=1$; Imhof's numerical inversion \citep{imhof1961}
is available when more accuracy is wanted. Every simulation in this paper uses the Satterthwaite
approximation. Neither computation needs a second model fit or resampling.

\subsection{The partition as a setting}\label{sec:partitionchoice}

\begin{theorem}[Growing partition]\label{thm:growingG}
Let the basis columns be the values at the group positions of fixed, continuously differentiable,
linearly independent functions on $[0,1]$, and let $G=G_n\le n$ with $G_n/\sqrt{\log n}\to\infty$.
Under the null hypothesis and conditions (A2), (A3) and (A4$'$) of Appendix~\ref{app:proofs}, $S$
converges to a weighted sum of $d$ chi-squared variables on one degree of freedom whose weights, given
in \eqref{eq:growingG}, do not depend on $G_n$. No lower bound on the number of observations per
group is required.
\end{theorem}

The rate condition is mild: the default partition below grows linearly in $n$ and satisfies it.

The Hosmer--Lemeshow statistic spends one degree of freedom per group, so a finer partition dilutes
it. The test statistic spends $d$ degrees of freedom whatever $G$ is, so the number of groups
becomes a setting that the analyst chooses, and we use two. The \emph{default partition} refines with
the sample, $G=\max(10,\lceil n/25\rceil)$, about twenty-five patients a group (the simulation study
rounded $n/25$ to the nearest integer, which moves the partition by at most one group); on clean data it gives
EDGE more power than ten groups (Section~\ref{sec:verdicts}). When records may be wrong,
the \emph{ten-group partition} keeps many clean records in the extreme groups, where corrupted records
collect, and protects the cubic basis much further (Section~\ref{sec:edgefr}); the symmetric basis is
better protected at the default partition against reversed records, though not against many
exaggerated ones (Section~\ref{sec:severity}). Pooling only the two extreme groups protects no better
than ten groups (Supporting Information S5).

In the R package \texttt{ebrahim.gof} \citep{ebrahim2025gof} the test is
\texttt{edge.gof(fit)} for a fitted \texttt{glm}. It reports two rows: the verdict at the default
partition and, as a robustness check, the same test at ten groups; its degrees of freedom are the
effective $\nu$ of the Satterthwaite reference. \texttt{G = 10} alone gives the ten-group partition, \texttt{basis = "sym"} or \texttt{"stukel"} for the other bases and
\texttt{weights = "score"} for the score form. For frozen predictions,
\texttt{edge.gof(y, predicted\_probs = p, external = TRUE, G = 10)} computes the external-mode
statistic at ten groups, with $\Omega=I_G$ and the constant column added, and reproduces every statistic
of Section~\ref{sec:realdata} at its own $G$; \texttt{run.all.external(y, p)} runs it at ten groups and at
the default partition beside the other external-validation tests of Supporting Information S7
(version 2.9.0 onwards).
